\section{One-sided squared sampling with one logarithm}
\label{sec:sampling}

The sampling estimate needed below concerns a squared one-sided loss.
It follows from a metric estimate in the proof of Lee's chaining theorem
\cite[Lemma~2.12, equation~(2.20)]{Lee2023}. We spell out the comparison:
a first-power ReLU sampling theorem alone would not give this result.
Throughout this section, $C$ denotes a universal constant that may change
from line to line, and $t_-:=\max\{-t,0\}$.

\begin{lemma}[A squared Lipschitz process]\label{lem:squared-process}
Let $a_1,\ldots,a_m\in\mathbb R^D$, and set
\[
 R=\max_{i\le m}\|a_i\|_2,
 \qquad V=\sup_{\|u\|_2\le1}
       \left(\sum_{i=1}^m\langle a_i,u\rangle^2\right)^{1/2}.
\]
For each $i$, let $\psi_i:\mathbb R\to\mathbb R$ satisfy
$\psi_i(0)=0$ and $\operatorname{Lip}(\psi_i)\le1$.
For independent uniform signs $\epsilon_i$,
\begin{equation}\label{eq:squared-process}
 \mathbb E_\epsilon\sup_{\|u\|_2\le1}
 \left|\sum_{i=1}^m\epsilon_i\psi_i(\langle a_i,u\rangle)^2\right|
 \le C R\sqrt{\log(2(m+D))}\,V.
\end{equation}
\end{lemma}

\begin{proof}
Use the genuine seminorms $N_i(u)=|\langle a_i,u\rangle|$ and the
symmetric quasimetric
\[
 d_A(u,v)^2
 =\sum_{i=1}^m N_i(u-v)^2\bigl(N_i(u)^2+N_i(v)^2\bigr).
\]
The stronger intermediate estimate established in the proof of
\cite[Lemma~2.12]{Lee2023} is
\begin{equation}\label{eq:lee-metric}
 \gamma_2(T,d_A)
 \le C\bigl(\lambda\sqrt{\log(2D)}+\kappa\bigr)
       \sup_{u\in T}\left(\sum_i N_i(u)^2\right)^{1/2},
 \qquad T\subseteq B_2^D,
\end{equation}
where
\[
 \lambda=\max_i\bigl(\mathbb E N_i(g)^2\bigr)^{1/2},
 \qquad \kappa=\mathbb E\max_i N_i(g),
 \qquad g\sim\mathcal N(0,I_D).
\]
Here $\gamma_2$ is the generic-chaining functional. This is equation
(2.20) of that proof, with the energy normalization (2.17) undone.
Seminorms follow from the norm statement by replacing
$N_i(u)$ with $(N_i(u)^2+\varepsilon^2\|u\|_2^2)^{1/2}$ and then
letting $\varepsilon\downarrow0$; the resulting metrics dominate
$d_A$, and the parameters on the right converge.
For our rank-one seminorms,
$\lambda=R$ and $\kappa\le R\sqrt{2\log(2m)}$.

Let $X_u=\sum_i\epsilon_i\psi_i(\langle a_i,u\rangle)^2$.
Its subgaussian increment metric is bounded by $\sqrt2\,d_A$, because
\begin{align*}
 &\sum_i\left(\psi_i(\langle a_i,u\rangle)^2
             -\psi_i(\langle a_i,v\rangle)^2\right)^2\\
 &\quad\le\sum_i |\langle a_i,u-v\rangle|^2
       \bigl(|\langle a_i,u\rangle|+|\langle a_i,v\rangle|\bigr)^2
 \le 2d_A(u,v)^2.
\end{align*}
We used both Lipschitz continuity and
$|\psi_i(t)|\le|t|$. The subgaussian chaining inequality, together
with $X_0=0$, bounds $\mathbb E\sup X_u$ by
$C\gamma_2(B_2^D,d_A)$. Apply the same argument to $-X_u$ and add.
Equation~\eqref{eq:lee-metric} now proves
\eqref{eq:squared-process}.
\end{proof}

\begin{remark}[The normalization in the metric estimate]
The proof of \cite[Lemma~2.12]{Lee2023} passes through a coarse additive
constant in its equation (2.26), after normalizing
$\sup_T\sum_iN_i^2=1$. The homogeneous version
\eqref{eq:lee-metric} follows by retaining the scale in the two geometric
terms of that argument. Indeed, in the rank-one case $T\subseteq B_2^D$
implies $N_i(u)\le\lambda$ on $T$; writing
$\|u\|_{\mathcal N}=\max_iN_i(u)$ gives
$d_A(u,v)\le\sqrt2\|u-v\|_{\mathcal N}$ on $T$.
Thus $\operatorname{diam}_{d_A}(T)\le C\lambda$ and the volumetric
packing constant in its Corollary~2.11 can be taken to be
$C\lambda$, since $T\subseteq\lambda B_{\mathcal N}$.
The additive term is therefore $C\lambda$, absorbed by
$C(\lambda\sqrt{\log(2D)}+\kappa)$.
The growth-functional portion of the proof is unchanged.
If the packing constant is less than one, first rescale the quasimetric
to apply that corollary in its stated normalization, and then scale back.
Most importantly, throughout that argument the $N_i$ are symmetric
seminorms. No assertion that a positive part is a seminorm is used.
\end{remark}

\begin{lemma}[Uniform one-sided squared sampling]
\label{lem:squared-sampling}
Let $(Z,Y)$ be a random pair with $Z\in\mathbb R^D$, $Y\in\{-1,1\}$,
$\mathbb E ZZ^{\mathsf T}=I_D$, and $\|Z\|_2^2=D$ almost surely.
For independent copies $(Z_i,Y_i)$, let
\[
 \Delta_n=\sup_{\|u\|_2\le1}
 \left|\frac1n\sum_{i=1}^n(Y_i\langle Z_i,u\rangle)_-^2
          -\mathbb E(Y\langle Z,u\rangle)_-^2\right|,
 \qquad h_n=\sqrt{\frac{D\log(2(n+D))}{n}}.
\]
Then $\mathbb E\Delta_n\le C h_n(1+C h_n)$.
In particular, for any fixed $\varepsilon,\delta>0$ there is a finite
$A_{\varepsilon,\delta}$ such that
\[
 n\ge A_{\varepsilon,\delta}D\log(2D)
 \quad\Longrightarrow\quad
 \mathbb P\{\Delta_n>\varepsilon\}\le\delta.
\]
The labels may depend arbitrarily on the features.
\end{lemma}

\begin{proof}
Put $Q_n=n^{-1}\sum_iZ_iZ_i^{\mathsf T}$.
Conditional on the observations, apply
Lemma~\ref{lem:squared-process} to
$a_i=Y_iZ_i/\sqrt n$ and $\psi_i(t)=t_-$.
Here $R=\sqrt{D/n}$ and $V=\sqrt{\|Q_n\|_{\mathrm{op}}}$.
Symmetrization gives
\begin{equation}\label{eq:sampling-symmetrization}
 \mathbb E\Delta_n
 \le C h_n\,\mathbb E\sqrt{\|Q_n\|_{\mathrm{op}}}.
\end{equation}
We also apply the same process lemma with $\psi_i(t)=t$.
Since $\mathbb E Q_n=I_D$, symmetrization and Jensen's inequality imply,
with $A=\mathbb E\|Q_n\|_{\mathrm{op}}$,
\[
 A\le1+\mathbb E\|Q_n-I_D\|_{\mathrm{op}}
 \le1+C h_n\mathbb E\sqrt{\|Q_n\|_{\mathrm{op}}}
 \le1+C h_n\sqrt A.
\]
Solving this quadratic inequality yields $\sqrt A\le1+C h_n$,
after increasing $C$. Substitute into
\eqref{eq:sampling-symmetrization}.
Finally, at $n\ge A_{\varepsilon,\delta}D\log(2D)$ the quantity
$h_n$ can be made uniformly small by increasing
$A_{\varepsilon,\delta}$. Markov's inequality proves the probability
statement. It is a statement for each specified sample size $n$.
\end{proof}

